\documentclass[aspectratio=169]{beamer}
\usetheme{Madrid}
\usecolortheme{whale}
\usepackage[utf8]{inputenc}
\usepackage[spanish]{babel}
\usepackage{amsmath,amssymb}
\usepackage{graphicx}
\usepackage{booktabs}
\usepackage{listings}
\usepackage{xcolor}

\definecolor{codegreen}{rgb}{0,0.6,0}
\definecolor{codegray}{rgb}{0.5,0.5,0.5}
\definecolor{codepurple}{rgb}{0.58,0,0.82}
\definecolor{backcolour}{rgb}{0.95,0.95,0.92}

\lstdefinestyle{pythonstyle}{
    backgroundcolor=\color{backcolour},
    commentstyle=\color{codegreen},
    keywordstyle=\color{blue},
    numberstyle=\tiny\color{codegray},
    stringstyle=\color{codepurple},
    basicstyle=\ttfamily\footnotesize,
    breakatwhitespace=false,
    breaklines=true,
    captionpos=b,
    keepspaces=true,
    showspaces=false,
    showstringspaces=false,
    showtabs=false,
    tabsize=2
}

\lstset{style=pythonstyle}

\title{Regresión Lineal Simple}
\subtitle{Inferencia y Predicción}
\author{Métodos de Análisis de Datos 1}
\institute{Universidad Nacional del Sur}
\date{Unidad 3 - Clase 15}

\begin{document}

\frame{\titlepage}

\begin{frame}{Contenidos}
\tableofcontents
\end{frame}

%==============================================================================
\section{Repaso Clase Anterior}
%==============================================================================

\begin{frame}{Repaso: Modelo y estimación}
\textbf{Modelo de regresión lineal simple:}
\[
Y_i = \beta_0 + \beta_1 X_i + \epsilon_i, \quad \epsilon_i \sim N(0, \sigma^2)
\]

\vspace{0.3cm}
\textbf{Estimadores por MCO:}
\[
\hat{\beta}_1 = \frac{S_{XY}}{S_{XX}}, \quad \hat{\beta}_0 = \bar{Y} - \hat{\beta}_1 \bar{X}
\]

\vspace{0.3cm}
\textbf{Bondad de ajuste:}
\[
R^2 = 1 - \frac{\text{SCR}}{\text{SCT}}
\]

\vspace{0.5cm}
\textbf{Hoy:} ¿Son significativos los coeficientes? ¿Cómo hacer predicciones con incertidumbre?
\end{frame}

%==============================================================================
\section{Distribución de los Estimadores}
%==============================================================================

\begin{frame}{Propiedades de $\hat{\beta}_1$ y $\hat{\beta}_0$}
Bajo los supuestos del modelo, los estimadores MCO siguen distribuciones normales:

\begin{block}{Pendiente}
\[
\hat{\beta}_1 \sim N\left(\beta_1, \frac{\sigma^2}{S_{XX}}\right)
\]
\end{block}

\begin{block}{Intercepto}
\[
\hat{\beta}_0 \sim N\left(\beta_0, \sigma^2 \left[\frac{1}{n} + \frac{\bar{X}^2}{S_{XX}}\right]\right)
\]
\end{block}

\vspace{0.3cm}
\textbf{Problema:} $\sigma^2$ es desconocido $\Rightarrow$ lo estimamos con:
\[
s^2 = \frac{\text{SCR}}{n - 2} = \frac{\sum (Y_i - \hat{Y}_i)^2}{n - 2}
\]
\end{frame}

\begin{frame}{Errores estándar}
Los \textbf{errores estándar} miden la precisión de las estimaciones:

\begin{block}{Error estándar de $\hat{\beta}_1$}
\[
\text{SE}(\hat{\beta}_1) = \sqrt{\frac{s^2}{S_{XX}}} = \frac{s}{\sqrt{S_{XX}}}
\]
\end{block}

\begin{block}{Error estándar de $\hat{\beta}_0$}
\[
\text{SE}(\hat{\beta}_0) = s \sqrt{\frac{1}{n} + \frac{\bar{X}^2}{S_{XX}}}
\]
\end{block}

\vspace{0.3cm}
\textbf{Nota:} Cuanto mayor sea $S_{XX}$ (mayor variabilidad en $X$), menor será el error estándar de $\hat{\beta}_1$
\end{frame}

%==============================================================================
\section{Tests de Hipótesis}
%==============================================================================

\begin{frame}{Test sobre la pendiente}
La pregunta más importante: \textbf{¿Hay relación entre $X$ e $Y$?}

\begin{block}{Hipótesis}
\[
H_0: \beta_1 = 0 \quad \text{vs} \quad H_1: \beta_1 \neq 0
\]
\end{block}

\textbf{Estadístico de prueba:}
\[
t = \frac{\hat{\beta}_1 - 0}{\text{SE}(\hat{\beta}_1)} = \frac{\hat{\beta}_1}{s/\sqrt{S_{XX}}} \sim t_{n-2}
\]

\textbf{Regla de decisión:} Al nivel $\alpha$, rechazamos $H_0$ si:
\[
|t| > t_{n-2, \alpha/2}
\]

\textbf{Valor-p:} $P(|T| > |t|)$ donde $T \sim t_{n-2}$
\end{frame}

\begin{frame}{Ejemplo: Test sobre pendiente}
Datos: Horas de estudio ($X$) vs Calificación ($Y$) de 20 estudiantes

\vspace{0.3cm}
Modelo ajustado: $\hat{Y} = 40 + 5.2 X$

\begin{itemize}
    \item $\hat{\beta}_1 = 5.2$
    \item $\text{SE}(\hat{\beta}_1) = 1.3$
    \item $t = \frac{5.2}{1.3} = 4.0$
\end{itemize}

\vspace{0.3cm}
Con $n = 20$, $t_{18, 0.025} = 2.101$

\vspace{0.3cm}
Como $|4.0| > 2.101$ $\Rightarrow$ \textbf{Rechazamos} $H_0$

\vspace{0.3cm}
\textbf{Conclusión:} Hay evidencia significativa ($\alpha = 0.05$) de que las horas de estudio afectan la calificación.
\end{frame}

\begin{frame}{Test sobre el intercepto}
Menos común, pero a veces relevante:

\begin{block}{Hipótesis}
\[
H_0: \beta_0 = 0 \quad \text{vs} \quad H_1: \beta_0 \neq 0
\]
\end{block}

\textbf{Estadístico:}
\[
t = \frac{\hat{\beta}_0}{\text{SE}(\hat{\beta}_0)} \sim t_{n-2}
\]

\vspace{0.5cm}
\textbf{¿Cuándo es útil?}
\begin{itemize}
    \item Cuando tiene sentido que la línea pase por el origen
    \item Ejemplo: relación dosis-respuesta donde dosis 0 $\Rightarrow$ respuesta 0
\end{itemize}
\end{frame}

%==============================================================================
\section{Intervalos de Confianza}
%==============================================================================

\begin{frame}{Intervalo de confianza para $\beta_1$}
Un intervalo de confianza del $(1-\alpha)\times 100\%$ para $\beta_1$ es:

\begin{block}{IC para la pendiente}
\[
\hat{\beta}_1 \pm t_{n-2, \alpha/2} \cdot \text{SE}(\hat{\beta}_1)
\]
\end{block}

\vspace{0.5cm}
\textbf{Interpretación:} Con 95\% de confianza, el verdadero valor de $\beta_1$ está en este intervalo.

\vspace{0.5cm}
\textbf{Relación con el test:} Si el IC no contiene 0, rechazamos $H_0: \beta_1 = 0$
\end{frame}

\begin{frame}{Ejemplo: IC para pendiente}
Continuando el ejemplo anterior:

\begin{itemize}
    \item $\hat{\beta}_1 = 5.2$
    \item $\text{SE}(\hat{\beta}_1) = 1.3$
    \item $t_{18, 0.025} = 2.101$
\end{itemize}

\vspace{0.3cm}
\textbf{IC 95\%:}
\[
5.2 \pm 2.101 \times 1.3 = 5.2 \pm 2.7 = [2.5, 7.9]
\]

\vspace{0.5cm}
\textbf{Interpretación:} Con 95\% de confianza, cada hora adicional de estudio aumenta la calificación entre 2.5 y 7.9 puntos.

\vspace{0.3cm}
Como el intervalo no contiene 0, confirmamos que $\beta_1 \neq 0$ (consistente con el test).
\end{frame}

\begin{frame}{Intervalo de confianza para $\beta_0$}
Análogamente:

\begin{block}{IC para el intercepto}
\[
\hat{\beta}_0 \pm t_{n-2, \alpha/2} \cdot \text{SE}(\hat{\beta}_0)
\]
\end{block}

\vspace{0.5cm}
\textbf{Nota:} Los IC para $\beta_0$ suelen ser muy anchos si $\bar{X}$ está lejos de 0, porque implica \textbf{extrapolación}.
\end{frame}

%==============================================================================
\section{ANOVA del Modelo de Regresión}
%==============================================================================

\begin{frame}{Tabla ANOVA para regresión}
Otra forma de probar $H_0: \beta_1 = 0$ es mediante ANOVA:

\begin{table}
\centering
\begin{tabular}{lccccc}
\toprule
\textbf{Fuente} & \textbf{SC} & \textbf{gl} & \textbf{CM} & \textbf{F} & \textbf{p-valor} \\
\midrule
Regresión & SCE & 1 & SCE/1 & CM$_E$/CM$_R$ & $P(F_{1,n-2} > F)$ \\
Residual & SCR & $n-2$ & SCR/$(n-2)$ & & \\
\midrule
Total & SCT & $n-1$ & & & \\
\bottomrule
\end{tabular}
\end{table}

\vspace{0.3cm}
donde:
\begin{itemize}
    \item SC = Suma de Cuadrados
    \item gl = grados de libertad
    \item CM = Cuadrado Medio = SC/gl
    \item $F = \frac{\text{CM}_E}{\text{CM}_R}$
\end{itemize}
\end{frame}

\begin{frame}{Relación entre test $t$ y test $F$}
En regresión lineal simple:

\begin{alertblock}{Equivalencia}
\[
F = t^2
\]
\end{alertblock}

\vspace{0.3cm}
Ambos tests prueban $H_0: \beta_1 = 0$ y dan la misma conclusión.

\vspace{0.5cm}
\textbf{¿Cuándo usar cada uno?}
\begin{itemize}
    \item Test $t$: regresión simple, o para coeficientes individuales en regresión múltiple
    \item Test $F$: regresión múltiple (próxima clase), para probar si \textbf{al menos un} predictor es útil
\end{itemize}
\end{frame}

%==============================================================================
\section{Predicción}
%==============================================================================

\begin{frame}{Dos tipos de predicción}
Dado un valor $X_0$, queremos estimar:

\begin{enumerate}
    \item \textbf{La media de $Y$ para $X = X_0$}
    \[
    E(Y | X = X_0) = \beta_0 + \beta_1 X_0
    \]
    $\Rightarrow$ \textbf{Intervalo de confianza}
    
    \vspace{0.5cm}
    \item \textbf{El valor de una nueva observación individual con $X = X_0$}
    \[
    Y_{\text{nuevo}} = \beta_0 + \beta_1 X_0 + \epsilon_{\text{nuevo}}
    \]
    $\Rightarrow$ \textbf{Intervalo de predicción}
\end{enumerate}
\end{frame}

\begin{frame}{Predicción puntual}
Para cualquiera de los dos casos, la estimación puntual es la misma:

\begin{block}{Predicción puntual}
\[
\hat{Y}_0 = \hat{\beta}_0 + \hat{\beta}_1 X_0
\]
\end{block}

\vspace{0.5cm}
La diferencia está en la \textbf{incertidumbre} asociada a la predicción.
\end{frame}

\begin{frame}{Intervalo de confianza para la media}
IC $(1-\alpha)\times 100\%$ para $E(Y|X=X_0)$:

\begin{block}{IC para la media}
\[
\hat{Y}_0 \pm t_{n-2, \alpha/2} \cdot s \sqrt{\frac{1}{n} + \frac{(X_0 - \bar{X})^2}{S_{XX}}}
\]
\end{block}

\vspace{0.3cm}
\textbf{Interpretación:} Estimamos la calificación \textbf{promedio} de todos los estudiantes que estudian $X_0$ horas.

\vspace{0.3cm}
\textbf{Nota:} El intervalo es más estrecho cerca de $\bar{X}$ y más ancho en los extremos.
\end{frame}

\begin{frame}{Intervalo de predicción para nueva observación}
IC $(1-\alpha)\times 100\%$ para $Y_{\text{nuevo}}$:

\begin{block}{Intervalo de predicción}
\[
\hat{Y}_0 \pm t_{n-2, \alpha/2} \cdot s \sqrt{1 + \frac{1}{n} + \frac{(X_0 - \bar{X})^2}{S_{XX}}}
\]
\end{block}

\vspace{0.3cm}
\textbf{Interpretación:} Estimamos la calificación de \textbf{un estudiante específico} que estudia $X_0$ horas.

\vspace{0.3cm}
\textbf{Diferencia clave:} El término adicional "$1$" dentro de la raíz representa la variabilidad individual $\epsilon$.
\end{frame}

\begin{frame}{Comparación de intervalos}
\begin{center}
\includegraphics[width=0.75\textwidth]{intervalos_prediccion.png}
\end{center}

\begin{itemize}
    \item \textcolor{blue}{Azul}: Intervalo de confianza para la media
    \item \textcolor{red}{Rojo}: Intervalo de predicción
    \item El de predicción es siempre más ancho
    \item Ambos son más anchos lejos de $\bar{X}$ (extrapolación riesgosa)
\end{itemize}
\end{frame}

\begin{frame}{Ejemplo numérico: Predicción}
Modelo: Calificación = $40 + 5.2 \times$ Horas

\begin{itemize}
    \item $n = 20$, $\bar{X} = 4$, $S_{XX} = 80$, $s = 8$
    \item Queremos predecir para $X_0 = 5$ horas
\end{itemize}

\vspace{0.3cm}
\textbf{Predicción puntual:} $\hat{Y}_0 = 40 + 5.2 \times 5 = 66$

\vspace{0.3cm}
\textbf{IC 95\% para la media:}
\[
66 \pm 2.101 \times 8 \sqrt{\frac{1}{20} + \frac{(5-4)^2}{80}} = 66 \pm 3.8 = [62.2, 69.8]
\]

\vspace{0.3cm}
\textbf{Intervalo de predicción 95\%:}
\[
66 \pm 2.101 \times 8 \sqrt{1 + \frac{1}{20} + \frac{(5-4)^2}{80}} = 66 \pm 17.1 = [48.9, 83.1]
\]
\end{frame}

%==============================================================================
\section{Implementación en Python}
%==============================================================================

\begin{frame}[fragile]{Inferencia en Python}
\begin{lstlisting}[language=Python]
import statsmodels.api as sm
from statsmodels.formula.api import ols

# Ajustar modelo
modelo = ols('Y ~ X', data=data).fit()

# Resumen completo (incluye tests t, p-valores, IC)
print(modelo.summary())

# Extraer coeficientes e ICs
print(modelo.params)  # coeficientes
print(modelo.bse)     # errores estandar
print(modelo.tvalues) # estadisticos t
print(modelo.pvalues) # p-valores

# IC 95% para coeficientes
print(modelo.conf_int(alpha=0.05))
\end{lstlisting}
\end{frame}

\begin{frame}[fragile]{Tabla ANOVA en Python}
\begin{lstlisting}[language=Python]
from statsmodels.stats.anova import anova_lm

# Tabla ANOVA
tabla_anova = anova_lm(modelo, typ=2)
print(tabla_anova)

# Output:
#           sum_sq     df         F    PR(>F)
# X        1234.5    1.0    16.00   0.0008
# Residual 1389.2   18.0       NaN      NaN

# F = 16.00, p-valor = 0.0008 < 0.05
# => Rechazamos H0: beta_1 = 0
\end{lstlisting}
\end{frame}

\begin{frame}[fragile]{Predicción en Python}
\begin{lstlisting}[language=Python]
import pandas as pd

# Nuevos datos
nuevos = pd.DataFrame({'X': [3, 5, 7]})

# Prediccion puntual
pred = modelo.predict(nuevos)
print("Predicciones:", pred.values)

# Con intervalos
pred_obj = modelo.get_prediction(nuevos)
intervalos = pred_obj.summary_frame(alpha=0.05)
print(intervalos)

# Columnas:
# - mean: prediccion puntual
# - mean_ci_lower/upper: IC para la media
# - obs_ci_lower/upper: intervalo de prediccion
\end{lstlisting}
\end{frame}

%==============================================================================
\section{Implementación en R}
%==============================================================================

\begin{frame}[fragile]{Inferencia en R}
\begin{lstlisting}[language=R]
# Ajustar modelo
modelo <- lm(Y ~ X, data = data)

# Resumen (incluye todo)
summary(modelo)

# Extraer componentes
coef(modelo)           # coeficientes
summary(modelo)$coef   # tabla completa con SE, t, p

# IC 95% para coeficientes
confint(modelo, level = 0.95)

# Tabla ANOVA
anova(modelo)
\end{lstlisting}
\end{frame}

\begin{frame}[fragile]{Predicción en R}
\begin{lstlisting}[language=R]
# Nuevos datos
nuevos <- data.frame(X = c(3, 5, 7))

# Prediccion puntual
predict(modelo, nuevos)

# IC para la media
predict(modelo, nuevos, interval = "confidence", 
        level = 0.95)

# Intervalo de prediccion
predict(modelo, nuevos, interval = "prediction", 
        level = 0.95)

# Output: data.frame con columnas fit, lwr, upr
\end{lstlisting}
\end{frame}

%==============================================================================
\section{Resumen}
%==============================================================================

\begin{frame}{Resumen de la clase}
\textbf{Aprendimos:}
\begin{itemize}
    \item Distribución de $\hat{\beta}_0$ y $\hat{\beta}_1$
    \item Tests de hipótesis: $H_0: \beta_1 = 0$
    \item Intervalos de confianza para los coeficientes
    \item Tabla ANOVA y relación $F = t^2$
    \item Dos tipos de predicción:
    \begin{itemize}
        \item IC para la media de $Y$
        \item Intervalo de predicción para nueva observación
    \end{itemize}
    \item Implementación completa en Python y R
\end{itemize}

\vspace{0.5cm}
\textbf{Próxima clase:} Diagnóstico del modelo (verificar supuestos, detectar problemas)
\end{frame}

\end{document}
